\cotefile{15}
\scanpage{170}

\uline{Foncteurs fibres motiviques, et \unsure{gpes} \unsure{fondamentaux} motiviques.}

À examiner :

0) VA \textcircled{à la Weil} \unsure{associées} aux motifs purs de poids impair \unsure{sous}
\unsure{Conditions} de Hodge--Riemann\ldots

\margin{OK *}
\textcircled{1} Extension des \unsure{groupes} \unsure{définis} par des \unsure{présentiques|présentations}
\uline{\unsure{Groupe} fondamental motivique.}

\textcircled{2} Relations avec le corps de classes

\margin{OK}
\textcircled{3} Cas du spectre d'un corps fini. Pb motivique
\unsure{points|partiel} à \unsure{val.} dans $\mathbb{R}$ ???

\margin{*}
4) Rôle des \unsure{gpes} fondamentaux motiviques pour la comparaison des réseaux canoniques
\diagram{\[
\left[
\begin{array}{ll}
\text{des Coh. de Hodge} &
\left|\begin{array}{l}
\text{2 réseaux \unsure{transcendts|transcendants} \unsure{pr.|pour} des places diff}\vphantom{\big|}\\
\text{réseau transcendant et réseau de Hodge}
\end{array}\right.\\[6pt]
\text{des Coh. $l$-adique} &\big|\ \text{2 réseaux transcendants}
\end{array}
\right]
\]}

5) Relations \inserted{de 4)} avec mes notes anciennes

6) Fonctions $L$ : équation fonctionnelle

\margin{OK}
7) Structures \unsure{positives}\ldots

\deleted{\unsure{7}}8) Topologie à mettre sur $\mathbb{Z}$ $-\, -\, -\, -$

9) Hyp. de Riemann ??

\margin{\unsure{Positivité} à la Hodge en se ramenant à des questions de \uline{diviseurs} par \unsure{éclatement|éclatements}, ou \unsure{vertu des|voie des} \unsure{pinceaux} ??}

\margin{169}
